\label{sec:discussion}

\subsection{The Algorithmic Indifference Standard}

We propose the following standard for evaluating whether an algorithmic redistricting
methodology is constitutionally sound on incumbency grounds:

\begin{definition}[Algorithmic Indifference Standard]
An algorithmic redistricting methodology satisfies the \emph{algorithmic indifference
standard} on incumbency if and only if:
\begin{enumerate}
  \item The algorithm receives no incumbency data (no incumbent home addresses, no
    prior district boundary data used for incumbency-protection purposes) as an input.
  \item The algorithm's output is determined solely by constitutionally required or
    permitted criteria (population balance, contiguity, compactness, community of
    interest) that do not include incumbency.
  \item The algorithm's incumbency outcomes are statistically consistent with what a
    random geographic redistricting process would produce given the same geographic
    constraints.
\end{enumerate}
\end{definition}

The \textsc{bisect} algorithm satisfies all three prongs. It receives no incumbent
addresses. Its optimisation target is population balance and geometric compactness.
And its incumbency outcomes (I.1--I.3) are statistically consistent with the geographic
random baseline across NC, WI, and TX.

A methodology that satisfies the algorithmic indifference standard cannot be challenged
on the ground that it ``fails to consider'' incumbency in any constitutionally meaningful
sense: the constitutional doctrine permits but does not require incumbency consideration,
and the algorithm's outputs are equivalent to what a truly neutral process would produce.

\subsection{Post-\textit{Rucho} State Court Applications}

The algorithmic indifference standard has particular relevance in state court proceedings
following \textit{Rucho v.\ Common Cause}\footnote{\textit{Rucho v.\ Common Cause},
588 U.S.\ 684 (2019).} State courts evaluating redistricting plans under state
constitutional free-elections, equal protection, or partisan fairness claims have
increasingly adopted the approach of comparing enacted maps to neutral algorithmic
baselines. The North Carolina Supreme Court in \textit{Harper v.\ Hall}\footnote{
\textit{Harper v.\ Hall}, 868 S.E.2d 499 (N.C.\ 2022).} and the Pennsylvania Supreme
Court in \textit{League of Women Voters v.\ Commonwealth}\footnote{\textit{League of
Women Voters v.\ Commonwealth}, 178 A.3d 737 (Pa.\ 2018).} both used neutral map
comparisons to evaluate whether enacted maps departed from what neutral redistricting
would produce.

In these proceedings, incumbents may argue that the algorithmic baseline is unfair
because it does not protect incumbency. The algorithmic indifference standard provides
the response: an algorithm that is indifferent to incumbency produces outcomes that
are statistically equivalent to what neutral redistricting---redistricting that is
directed to achieve only constitutionally required criteria---would produce. The
algorithm's outcomes are, by construction, the outputs of a fair process that applies
no thumb on the scale for or against any specific incumbent.

\subsection{The Expert Witness Application}

For expert witnesses proposing \textsc{bisect} maps in redistricting litigation, the
legal analysis in this paper supports three specific arguments.

\paragraph{Incumbency consideration is not required.}
The first and most direct argument: federal constitutional doctrine does not require
any redistricting methodology to consider incumbency. \textit{Karcher}'s identification
of incumbency as a legitimate objective means that mapmakers may consider it; it does
not mean that methodologies that do not consider it are constitutionally defective.
An algorithm that ignores incumbency is no different in constitutional status from a
commission that is instructed not to consider incumbent addresses in drawing its map---
a perfectly lawful instruction that courts have endorsed in Arizona's and California's
commission systems.

\paragraph{Algorithmic outcomes are equivalent to neutral redistricting.}
The second argument draws on I.1--I.3: the \textsc{bisect} algorithm's incumbency
outcomes are statistically indistinguishable from what a random geographic redistricting
process would produce. This means that the algorithm is not more disruptive to incumbency
than any other neutral redistricting process. An expert witness can testify that the
algorithmic map's pairing rate, safe-seat fraction, and open-seat count are all within
the expected range for a geographically constrained redistricting process with no
incumbency input---the algorithm is neutral on incumbency because its outputs reflect
geographic chance, not strategic choice.

\paragraph{The enacted map's incumbency outcomes are non-random.}
The third argument is the flip side: I.1's binomial tests demonstrate that the enacted
map's incumbency outcomes---zero pairings in NC and WI, significantly below-baseline
pairing in TX---are inconsistent with random redistricting. This is circumstantial
evidence that the enacted mapmakers used incumbent address data to engineer protective
districts. The contrast between algorithmic neutrality and enacted strategic protection
provides courts with a precise quantitative measure of the incumbency manipulation
that the enacted map embeds.

\subsection{Practitioner Decision Table}

Table~\ref{tab:decision-table} summarises the legal analysis for practitioners
evaluating whether to deploy \textsc{bisect} maps in redistricting proceedings.

\begin{table}[ht]
\centering
\begin{threeparttable}
\caption{Decision table for incumbency objections to algorithmic redistricting.}
\label{tab:decision-table}
\begin{tabular}{p{3.5cm} p{4cm} p{5cm}}
\toprule
Objection & Legal Standard & Response \\
\midrule
Algorithm ignores incumbency & \textit{Karcher}: permissive, not mandatory & No court
  requires incumbency consideration; algorithm satisfies all required criteria \\
\addlinespace
Algorithm produces pairing & No constitutional prohibition on pairing & Pairing rate
  is consistent with geographic baseline; enacted zero-pairing requires incumbency data \\
\addlinespace
Algorithm produces unsafe seats & VRA requires minority opportunity, not safe seats &
  VRA compliance is evaluated ex post; majority-minority districts preserved where required \\
\addlinespace
State statute permits incumbency & Permit $\neq$ require & Algorithm is not prohibited
  from ignoring incumbency; no state statute mandates protection \\
\addlinespace
Algorithm disrupts continuity & \textit{Burdick}: election admin interest & Algorithm
  disruption is equivalent to neutral redistricting; more continuity than random \\
\bottomrule
\end{tabular}
\end{threeparttable}
\end{table}
